\chapter{ภาคผนวก ข — ปัญญาประดิษฐ์ โค้ดคอมพิวเตอร์ และชุดคำสั่งพรอมพ์สำหรับการเกษตร}
\label{appendix:ai_code}

\section{โค้ดภาษา Python สำหรับพยากรณ์ประสิทธิภาพจุลินทรีย์ดินด้วย Random Forest}

สคริปต์นี้สาธิตการใช้แบบจำลองการเรียนรู้ของเครื่อง (Machine Learning: Random Forest Regressor) เพื่อทำนายปริมาณชีวมวลจุลินทรีย์ดินและอัตราการปลดปล่อยฟอสเฟตจากข้อมูลเซนเซอร์โทรมาตร (IoT Soil Telemetry)

\begin{lstlisting}[language=Python, caption={การพยากรณ์ฟอสฟอรัสละลายได้จากแบคทีเรียละลายฟอสเฟตด้วย Scikit-Learn}]
import numpy as np
import pandas as pd
from sklearn.ensemble import RandomForestRegressor
from sklearn.model_selection import train_test_split
from sklearn.metrics import mean_squared_error, r2_score

# 1. จำลองชุดข้อมูลเซนเซอร์ดินและค่าฟอสฟอรัสละลายได้ (Available P)
np.random.seed(42)
n_samples = 300

data = pd.DataFrame({
    'soil_temp_c': np.random.uniform(22, 36, n_samples),
    'soil_moisture_pct': np.random.uniform(15, 65, n_samples),
    'soil_ph': np.random.uniform(5.0, 7.8, n_samples),
    'soil_ec_uscm': np.random.uniform(150, 950, n_samples),
    'pgpr_inoculum_density_log': np.random.uniform(5.5, 8.5, n_samples)
})

# ค่าเป้าหมาย: อัตราฟอสฟอรัสที่ละลายได้ (mg/kg) ขึ้นกับ pH ความชื้น และปริมาณหัวเชื้อ
data['available_p_mg_kg'] = (
    12.5 
    + 3.2 * data['pgpr_inoculum_density_log']
    + 0.45 * data['soil_moisture_pct']
    - 2.1 * (data['soil_ph'] - 6.5)**2
    + np.random.normal(0, 1.2, n_samples)
)

# 2. แบ่งชุดข้อมูลฝึกสอนและทดสอบ
X = data[['soil_temp_c', 'soil_moisture_pct', 'soil_ph', 'soil_ec_uscm', 'pgpr_inoculum_density_log']]
y = data['available_p_mg_kg']

X_train, X_test, y_train, y_test = train_test_split(X, y, test_size=0.2, random_state=42)

# 3. สร้างและฝึกสอนแบบจำลอง Random Forest
rf_model = RandomForestRegressor(n_estimators=100, max_depth=6, random_state=42)
rf_model.fit(X_train, y_train)

# 4. ประเมินผลแบบจำลอง
y_pred = rf_model.predict(X_test)
print(f"R-squared Score: {r2_score(y_test, y_pred):.3f}")
print(f"Root Mean Squared Error (RMSE): {np.sqrt(mean_squared_error(y_test, y_pred)):.3f} mg/kg")

# 5. ฟังก์ชันทำนายค่าสำหรับระบบ IoT
def predict_soil_p(temp, moisture, ph, ec, inoculum_log):
    sample = np.array([[temp, moisture, ph, ec, inoculum_log]])
    return rf_model.predict(sample)[0]

# ตัวอย่างการทำนาย: ดินอุณหภูมิ 28C, ความชื้น 45%, pH 6.3, EC 400, หัวเชื้อ 10^7 cfu/g
pred_p = predict_soil_p(28.0, 45.0, 6.3, 400.0, 7.0)
print(f"คาดการณ์ฟอสฟอรัสละลายได้: {pred_p:.2f} mg/kg")
\end{lstlisting}

\section{โค้ดภาษา Python สำหรับตรวจจับรอยโรคพืชด้วยสถาปัตยกรรม Edge AI}

\begin{lstlisting}[language=Python, caption={การประมวลผลและเตรียมภาพใบพืชเพื่อจำแนกโรคจากเชื้อราบนอุปกรณ์พกพา}]
import cv2
import numpy as np

def preprocess_leaf_roi(image_path, target_size=(224, 224)):
    img = cv2.imread(image_path)
    # แปลงสีเป็น HSV เพื่อตัดพื้นหลังดินและอุปกรณ์ออก
    hsv = cv2.cvtColor(img, cv2.COLOR_BGR2HSV)
    # กรองเฉพาะช่วงสีเขียวและสีเหลืองของใบพืช
    lower_green = np.array([25, 40, 40])
    upper_green = np.array([85, 255, 255])
    mask = cv2.inRange(hsv, lower_green, upper_green)
    
    # ปรับปรุงขอบเขตด้วย Morphological Operations
    kernel = np.ones((5, 5), np.uint8)
    mask = cv2.morphologyEx(mask, cv2.MORPH_CLOSE, kernel)
    
    # หาพื้นที่รอบนอกของใบพืช
    contours, _ = cv2.findContours(mask, cv2.RETR_EXTERNAL, cv2.CHAIN_APPROX_SIMPLE)
    if contours:
        c = max(contours, key=cv2.contourArea)
        x, y, w, h = cv2.boundingRect(c)
        leaf_crop = img[y:y+h, x:x+w]
    else:
        leaf_crop = img
        
    resized = cv2.resize(leaf_crop, target_size)
    normalized = resized.astype(np.float32) / 255.0
    return np.expand_dims(normalized, axis=0)
\end{lstlisting}

\section{ชุดคำสั่งพรอมพ์ (System Prompts) ด้านจุลชีววิทยาการเกษตร}

\begin{keyconcept}[System Prompt: ผู้เชี่ยวชาญการวินิจฉัยโรคพืชและออกแบบแผนชีวภัณฑ์]
\textbf{Role:} คุณคือนักวิชาการเกษตรผู้เชี่ยวชาญด้านโรคพืชและจุลชีววิทยาควบคุมชีววิธี (Plant Pathologist \& Biocontrol Specialist)

\textbf{Task:} วินิจฉัยเชื้อก่อโรคจากอาการที่ปรากฏบนพืชเศรษฐกิจ (เช่น ทุเรียน, ข้าว, มันสำปะหลัง) และออกแบบแผนการใช้ชีวภัณฑ์ทดแทนสารเคมีอย่างปลอดภัย

\textbf{Input Details:} ชนิดพืช, ชิ้นส่วนที่เกิดโรค, ลักษณะรอยแผล (ช้ำน้ำ, เส้นใยราขาว, สปอร์ดำ, เมือกแบคทีเรีย), สภาพอากาศช่วงที่ระบาด

\textbf{Output Requirements:}
\begin{enumerate}
    \item ระบุชนิดเชื้อก่อโรคที่เป็นไปได้ (\textit{Genus species}) พร้อมระดับความน่าจะเป็น
    \item แนะนำชนิดจุลินทรีย์ปฏิปักษ์ที่เหมาะสม (เช่น \organism{Trichoderma harzianum}, \organism{Bacillus subtilis})
    \item อัตราส่วนและเวลาการพ่นชีวภัณฑ์ที่ให้ประสิทธิภาพสูงสุด (หลีกเลี่ยงแสงแดดจัดและสารเคมีกำจัดเชื้อรา)
\end{enumerate}
\end{keyconcept}
